\infowidth{8cm} % 设置封面信息下划线的宽度，默认为8cm

\cheading{天津大学2026届本科生毕业论文}       % 设置正文的页眉，需要填上对应的毕业年份
\ctitle{大学生创新创业教育模式对其创业意愿的影响机制研究}       % 封面用论文标题
\caffil{管理与经济学部}                    % 学院名称
\csubject{工商管理}                        % 专业名称
\cgrade{2022级}                            % 年级
\cauthor{杨绚伊}                           % 学生姓名
\cnumber{3022209135}                       % 学生学号
\csupervisor{肖应钊}                       % 导师姓名

\cabstractcn{
    在国家大力推动\nolinebreak[4]“大众创业、万众创新”的大环境之下，高校创新创业教育课程发展趋向多样化。现有的研究还没有对不同的课程模式对学生的创业意愿产生怎样的具体影响途径进行系统阐述。本研究以自我决定理论、计划行为理论为依托，提出了课程模式、自主动机、创业意愿之间的中介效应模型，并以天津某高校创业学院154名参加创新创业课程学习的学生为研究对象，使用问卷调查法和SPSSAU软件进行了实证分析。研究结果表明，三种类型的课程都提高了学生自主动机，自主动机在各个类型的课程和创业意愿之间起完全中介作用，外部动机只在融合型课程中对两者的关系产生负向调节作用。该发现显示创新创业教育依靠增进学生自主性从而激发其创业意向的基本机理，给高校优化课程结构来调动学生内在潜能赋予了重要的理论支撑和实际参照依据。
}

\ckeywordcn{
    创新创业教育模式，自主动机，外部动机，创业意愿，中介模型
}

\cabstracten{
    Under the environment of the country's vigorous promotion of ``mass entrepreneurship and innovation,'' the development of innovation and entrepreneurship education courses in colleges and universities tends to diversify. The existing research has not systematically expounded the specific ways of how different curriculum models affect students' entrepreneurial intention. Based on the theory of self-determination and the theory of planned behavior, this thesis puts forward the mediating effect model between curriculum model, autonomous motivation and entrepreneurial intention, and takes 154 students who participate in the course of innovation and entrepreneurship in an entrepreneurship college in Tianjin as the research object. The questionnaire survey method and SPSSAU software were used for empirical analysis. The results show that all three types of courses have improved students' autonomous motivation. Autonomous motivation plays a complete mediating role between various types of courses and entrepreneurial intention, and external motivation only has a negative moderating effect on the relationship between the two in the integrated courses. This finding shows that innovation and entrepreneurship education relies on the basic mechanism of enhancing students' autonomy to stimulate their entrepreneurial intention, which gives important theoretical support and practical reference for colleges and universities to optimize the curriculum structure to mobilize students' intrinsic potential.
}

\ckeyworden{
    Innovation and Entrepreneurship Education Model；Autonomous Motivation；External Motivation；Entrepreneurial Intention；Moderated Mediation Model
}
